% Last updated: 2025/04/24

% AND
Surgery is one of the most expensive categories of healthcare, estimated to represent about 50\% of all Medicare costs in 2014 (\cite{kaye_understanding_2020}). In recent years, policymakers have increased efforts to decrease healthcare expenditures through programs that target hospital-level spending, with less direct focus on those who actually make treatment decisions: physicians. Surgery is often not the only treatment option, but a physician may still choose to operate in cases despite low clinical benefit. Spine surgery is one such example. According to the Lown Institute, these instances of low-value care make up 11-14\% of some spine surgeries in the Medicare population, adding up to about \$2 billion dollars from 2019 to 2021 (\cite{noauthor_unnecessary_nodate}). 

% BUT
An individual physician's choice the type of operation and whether to operate in the first place plays a considerable role in surgery costs. These tendencies are facets of their practice style, which is often quantified in the literature as the the propensity of intensive treatment. The formation of practice styles begins in medical school, continues through additional training, namely residency and fellowship, and even continues to evolve in a physician's current practice environment. Previous literature traces the influence of medical school and residency on current practice style through program rankings, which serves as a proxy for quality of training (\cite{epstein2009}; \cite{schnell_addressing_2018}). The implicit assumption built into these papers is that training plays a large role in shaping a physician's beliefs about the standard of care, and that physicians practice similarly to how they are taught as a result. However, papers that only consider program rankings are unable to fully comment on how the material and knowledge passed down through training directly shapes a physician's current practice. 

% THEREFORE
To address the question of how training impacts current practice style, I create a new measure of the practice intensity around a medical school that quantifies the standard of care that is taught at that medical school. Using this measure, I analyze the effect of medical school on a physician's practice intensity post-training using a sample of newly practicing surgeons alongside a physician mover design. By comparing surgeons who went to medical school in different locations and controlling for current practice market effects, I am able to identify the effect of medical school.

% Discuss results
In my main analysis, I find no evidence that medical school has a statistically significant effect on a surgeon's practice style. In particular, a 1 percentage point increase in medical school intensity corresponds to a 0.08\% decrease in current intensity. These findings suggest that medical schools likely convey similar foundational knowledge about the types of patients suitable for operation to all medical students, but the application of that knowledge to actual practice intensity is refined in later stages of training after medical school. The perpetuation of differences in treatment intensity and healthcare expenditures is likely not due to differences in medical school teaching across the US. 